% ============================================================================
%   RÉSUMÉ (Français) — Aivancity: 250-300 mots, ≥5 mots-clés
% ============================================================================

\chapter*{Résumé}
\addcontentsline{toc}{chapter}{Résumé}
\markboth{Résumé}{}

L'appariement patient--essai clinique est un goulot d'étranglement de la 
recherche médicale, et les systèmes récents fondés sur les grands modèles 
de langage (LLM) tels que TrialGPT visent à automatiser l'évaluation de 
l'éligibilité des patients face aux critères d'inclusion et d'exclusion. 
Leur fiabilité probabiliste n'a cependant pas fait l'objet d'un audit 
systématique : une revue récente de vingt-quatre systèmes n'en identifie 
aucun qui rapporte des métriques de calibration --- alors qu'un modèle 
sur-confiant qui se trompe avec certitude est cliniquement dangereux.

Nous avons mené le premier audit systématique de la calibration des LLM 
open source dans ce domaine : douze configurations (six modèles $\times$ 
deux prompts) évaluées sur les 1{,}015 paires patient--critère de 
TrialGPT-Criterion-Annotations, par scoring direct des log-probabilités 
et quatre hypothèses pré-enregistrées (permutation, Bonferroni).

Les douze configurations dépassent le seuil clinique de $0{,}05$ (ECE de 
$0{,}124$ à $0{,}602$), et le RLHF dégrade significativement la 
calibration sur les deux familles ($p < 0{,}003$). Cinq configurations 
présentent un effondrement de mode, et les deux LLM médicaux atteignent 
une exactitude $\leq 0{,}4\%$ sur les $238$ cas discriminants. Cet 
effondrement est cependant largement un artefact de décodage : le signal 
discriminant subsiste dans les log-probabilités (AUC jusqu'à $0{,}91$) 
mais est masqué par un a priori de réponse. Un scalaire unique, le ratio 
signal/a priori, prédit l'effondrement, explique pourquoi une 
recalibration monotone ne peut le corriger, et justifie un re-seuillage 
qui restaure jusqu'à $+0{,}79$ de macro-F1. L'analyse distingue enfin 
l'effondrement récupérable de la corruption irrécupérable, distinction 
qui se réplique sur le corpus indépendant TREC Clinical Trials.

Les métriques agrégées peuvent ainsi masquer un échec complet sur les 
cas décisifs comme un signal récupérable sous une défaillance apparente. 
Un protocole d'évaluation responsable doit combiner calibration, 
diversité prédictive, analyse stratifiée par difficulté, et un diagnostic 
distinguant une règle de décision effondrée d'un signal corrompu.

\vspace{0.5cm}
\noindent\textbf{Mots-clés :} grands modèles de langage, calibration, 
appariement d'essais cliniques, RLHF, IA médicale, mode collapse, 
expected calibration error, re-seuillage.

\cleardoublepage